\section{Related Work}
\label{sec:related work}

Modern recommendation systems are based on Deep Learning Recommendation Models (DLRMs) and how to effectively model the feature interactions is a crucial factor for DLRMs~\cite{cheng2016wide,guo2017deepfm,huang2022neural,qu2018product,wang2024flip,wang2021dcn,zhang2022dhen}. Wide\&Deep~\cite{cheng2016wide} is one of the earliest efforts, which combines the logistic regression (wide part) and DNN (deep part) to capture low-order and high-order feature interactions respectively. DeepFM~\cite{guo2017deepfm} is another achievement, which integrates the Factorization Machine (FM) and the DNN. In addition, DeepCross~\cite{shan2016deep} is an extension of the residual network~\cite{he2016deep}, aiming to learn automatic feature interactions implicitly. While it has proven to be quite challenging to simply rely on DNNs to learn high-order feature interactions~\cite{qu2018product,rendle2020neural}.
Explicit cross methods design different operators to explicitly capture high-order feature interactions, such as PNN~\cite{qu2018product}, DCN~\cite{wang2017deep} and its successor DCNv2~\cite{wang2021dcn}, xDeepFM~\cite{lian2018xdeepfm}, FGCNN~\cite{liu2019feature} and FiGNN~\cite{li2019fi}. AutoInt~\cite{song2019autoint} and Hiformer~\cite{gui2023hiformer} employ attention mechanism with residual connections to learn complex interactions. DHEN~\cite{zhang2022dhen} proposes to combine multiple interaction operators together. Despite improving accuracy, these new architectures increase the model latency and memory consumption and their model size is relatively small.


Scaling law has emerged as a fundamental theme in deep learning and a pivotal catalyst for numerous breakthroughs over the past decade, particularly in Natural Language Processing (NLP)~\cite{kaplan2020scaling,hoffmann2022training}, Computer Vision (CV)~\cite{dosovitskiy2020image,zhai2022scaling}, and multi-modality modeling~\cite{radford2021learning,ramesh2021zero}, which describe the power-law correlations between model performance and scaling factors, such as model size, data volume, and computational capacity~\cite{kaplan2020scaling,achiam2023gpt,touvron2023llama,team2024gemma}.
Recently, scaling laws in recommendation systems have drawn much attention from researchers~\cite{guo2024scaling}. Studies have explored the scaling strategy for pre-training user activity sequences~\cite{chitlangia2023scaling}, general-purpose user representations~\cite{shin2023scaling,zhang2024scaling} and online retrieval~\cite{fang2024scaling,wang2024scaling}. Wukong~\cite{zhangwukong} stacks FMs and LCB to learn the feature interactions. Orthogonally, \citet{zhang2024scaling} scaled up a sequential recommendation model to 0.8B parameters. HSTU~\cite{zhaiactions} enhances the scaling effect of Generative Recommenders (GRs) which focus more on the sequence part.
